\documentclass[12pt]{article}
\usepackage{amsmath,amssymb,amsfonts}
\usepackage[margin=1in]{geometry}

\begin{document}

\section*{Chapter 2, Problem 13}

\subsection*{Problem Statement}
Prove that a particle moving under gravity in a plane from a fixed point $P$ to a vertical line $L$ will reach the line in minimum time by following the cycloid from $P$ to $L$ that intersects $L$ at right angles.

\subsection*{Solution}

This is a brachistochrone problem with one fixed endpoint and one free endpoint (the particle must reach the vertical line $L$, but the exact point on $L$ is free).

Set up coordinates with the $y$-axis pointing downward and $P$ at the origin. The line $L$ is the vertical line $x = a$ for some constant $a > 0$. The particle starts from rest, so its speed at depth $y$ is $v = \sqrt{2gy}$.

The travel time is
\[
  T = \int_0^{x_f} \frac{\sqrt{1+y'^2}}{\sqrt{2gy}}\,dx,
\]
where $y' = dy/dx$ and the upper limit $x_f$ corresponds to the point where the path meets $L$.

\subsubsection*{Step 1: The extremizing path is a cycloid}

From the standard brachistochrone analysis (Section 2.1--2.3 of the text), the Euler--Lagrange equation for $f = \sqrt{1+y'^2}/\sqrt{2gy}$ yields the cycloid:
\[
  x = a(\phi - \sin\phi), \qquad y = a(1-\cos\phi),
\]
where $a$ is a parameter (the radius of the generating circle) and $\phi$ is the cycloid parameter.

\subsubsection*{Step 2: Transversality condition at the free endpoint}

When one endpoint is free to move along a curve (here the vertical line $x = a$), the natural boundary condition (transversality condition) must be satisfied. For a functional $I = \int f(x,y,y')\,dx$ with the right endpoint free to slide along the curve $x = \text{const}$ (a vertical line), the transversality condition is:
\[
  \left.\frac{\partial f}{\partial y'}\right|_{x=a} \cdot \delta y + \left[f - y'\frac{\partial f}{\partial y'}\right]_{x=a} \cdot \delta x = 0.
\]

Since the endpoint moves along $x = a$, we have $\delta x = 0$ and $\delta y$ is arbitrary. Therefore:
\[
  \left.\frac{\partial f}{\partial y'}\right|_{x=a} = 0.
\]

Computing:
\[
  \frac{\partial f}{\partial y'} = \frac{y'}{\sqrt{2gy}\sqrt{1+y'^2}}.
\]

Setting this to zero requires $y' = 0$ at the endpoint on $L$, i.e., $dy/dx = 0$.

\subsubsection*{Step 3: Geometric interpretation}

The condition $y' = 0$ (i.e., $dy/dx = 0$) at the line $x = a$ means the curve has a horizontal tangent where it meets $L$. Since $L$ is vertical, a horizontal tangent is perpendicular to $L$.

Therefore, the cycloid meets the vertical line $L$ at right angles.

\subsubsection*{Step 4: Verification on the cycloid}

On the cycloid, $dx/d\phi = a(1-\cos\phi)$ and $dy/d\phi = a\sin\phi$, so:
\[
  y' = \frac{dy}{dx} = \frac{\sin\phi}{1-\cos\phi} = \cot\frac{\phi}{2}.
\]
The condition $y' = 0$ gives $\cot(\phi/2) = 0$, i.e., $\phi/2 = \pi/2$, so $\phi = \pi$. At this point on the cycloid, the tangent is horizontal, confirming the perpendicular intersection with the vertical line $L$.

\medskip
\noindent\textbf{Conclusion:} The minimum-time path from $P$ to the vertical line $L$ is the cycloid that intersects $L$ at right angles (with a horizontal tangent at the intersection point).

\end{document}
